% Intro -- paper plan Task 3.
% Rule: no measured numbers here (they live in Results via \input/macros).

\section{Introduction}

\subsection{Motivation: Why LLMs Hallucinate}
Modern Large Language Models (LLMs) demonstrate remarkable fluency across diverse natural language tasks. However, beneath this surface coherence lies a fundamental vulnerability: LLMs are trained to maximize next-token statistical likelihood, not empirical factuality. Because the training objective prioritizes token co-occurrence over logical validity, models frequently generate statements that are syntactically plausible yet factually false---a phenomenon known as \emph{hallucination}. Benchmark evaluations demonstrate that even frontier models emit unsupported atomic claims at rates that preclude their deployment in high-stakes domains such as healthcare, legal reasoning, and scientific question answering \cite{ravichander2025halogen}.

To mitigate hallucinations, practitioners have developed several paradigms, each focusing on a distinct stage of generation:
\begin{enumerate}
  \item \textbf{Pre-Generation Retrieval (Standard RAG):} Systems retrieve external passages before answering \cite{asai2024selfrag,yan2024crag}. While RAG grounds answers in corpus text, it typically queries unconditionally for every prompt, introducing substantial retrieval latency. Furthermore, models can still misinterpret retrieved passages or inject unsupported assertions alongside grounded text.
  \item \textbf{Post-Generation Verification:} Systems decompose emitted answers into atomic claims and verify them against reference sources \cite{min2023factscore,wei2024longformfactuality}. However, verification alone does not repair the error.
  \item \textbf{Monolithic Revision:} Methods such as Chain-of-Verification (CoVe) \cite{dhuliawala2023cove} or Self-Refine \cite{madaan2023selfrefine} detect errors and prompt the model to rewrite the entire response. Rewriting monolithic text is computationally wasteful and frequently introduces new hallucinations or erases previously accurate sentences.
\end{enumerate}

What has been missing is an inference architecture that unifies these stages around a \emph{shared, persistent claim state}: a dynamic representation that tracks each atomic statement, its evidence sufficiency, and its factual support, enabling surgical, sentence-level repair without full-text regeneration.

\subsection{The HealGen Approach}
We present \textbf{HealGen}, a claim-level, evidence-sufficiency-guided self-healing inference architecture designed to be accessible, efficient, and transparent for students and practitioners alike. HealGen operates on three core principles:
\begin{itemize}
  \item \textbf{Draft-First Decomposition:} The model first generates an unconstrained draft. The draft is immediately parsed into atomic, self-contained claims. This ensures verification operates on what the model actually asserted rather than speculative query hypotheses.
  \item \textbf{Sufficiency-Gated Retrieval:} Instead of retrieving blindly for every prompt, HealGen queries external evidence only when an Adaptive Evidence Sufficiency Predictor (AESP) determines that available evidence is inadequate ($\mathrm{score} < \tau$). This isolates retrieval effort strictly to uncertain or unsupported claims.
  \item \textbf{Sentence-Localized Splice:} Unsupported claims are tracked in a bipartite Dynamic Claim--Evidence Graph (DCEG). Rather than regenerating the entire response, HealGen repairs \emph{only the specific sentences} owning unsupported claims, preserving verified sentences 100\% byte-identical.
\end{itemize}

\subsection{Contributions}
This paper provides an empirical architecture and measurement report:
\begin{enumerate}
  \item \textbf{Architecture.} A closed, claim-level control loop combining draft-first decomposition, sufficiency-gated retrieval, bipartite evidence graphs, and sentence-localized repair.
  \item \textbf{Rigorous Benchmark Evaluation.} An open, fully local prototype (CPU-scale, no API keys) evaluated with a paired experimental design on HaluEval-QA against vanilla generation, always-on RAG, and CoVe under a frozen shared-pool protocol.
  \item \textbf{Calibration \& Ablations.} Fine-grained claim-level calibration diagnostics (reliability diagram, ECE, Brier score) and ablation experiments isolating the precise value of risk estimation, sufficiency gating, and localized repair.
\end{enumerate}

\subsection{Paper Roadmap for Learners}
To facilitate structured learning, this paper is organized as follows:
Section~\ref{sec:related} surveys related literature across retrieval, verification, and editing.
Section~\ref{sec:architecture} presents the HealGen architecture with a step-by-step pedagogical walkthrough of a concrete query.
Section~\ref{sec:method} provides formal mathematical and algorithmic formulations for all modules.
Section~\ref{sec:experiments} details our reproducible experimental protocol.
Section~\ref{sec:results} analyzes empirical findings, cost trade-offs, and calibration.
Section~\ref{sec:limitations} discusses practical engineering limitations, and Section~\ref{sec:conclusion} concludes with educational and architectural takeaways.
